%!TEX root = ../thesis.tex

\addchap*{Abstract}
\markboth{Abstract}{}

Coverage-guided greybox fuzzing is an important tool in ensuring the security of software.
It executes millions of inputs and keeps the ones that reach new code without
forming an explicit model of what inputs are good and how a target program behaves,
relying on raw throughput instead.
Which inputs to use, which inputs to modify and how to modify them is mostly
left to hand-tuned heuristics and does not adapt to the program being tested.

We frame emulator-based greybox fuzzing as a Markov decision process and apply EfficientZero
to it, a sample-efficient MuZero variant that plans with Monte-Carlo tree search inside a
learned latent model of its environment.
To the best of our knowledge, this is the first fuzzer that plans with a learned model of the target's behavior.
Prior RL-based fuzzers usually use model-free variants of Q-learning, and planning-based fuzzers
apply Monte-Carlo tree search to an explicitly constructed model.
Basic-block coverage collected from an emulator is used as an observation as well as to derive a reward.
We formulate two action spaces: input-level actions, where an action is an input the program
consumes, and corpus-level actions, where an action selects a position in a seed and the byte value
to write to it.
This approach is implemented in the fuzzer \emph{mugiwara} using the Rust programming language
and the MOGI emulation framework.

We evaluate six targets, three in the input-action setting and three in the corpus-action setting.
The agent solves all three input-action targets: it recovers a 16-byte passcode one position at a
time and activates an OPC~UA session against open62541.
On the corpus-action targets it beats the random control baseline at identical execution budget
and improves on it by a factor of \num{1.77} in the best run.
It also manages to derive structure related to the target program's input grammar:
for the cJSON and libxml2 parsers, it generates easily recognizable nested structure.
This serves to prove the feasibility of our fuzzer.

Limitations of the approach include throughput and stability.
Guest execution rates of \num{21} to \num{23} inputs per second fall three orders of magnitude short
of native coverage-guided fuzzing, and training is shown to be unstable and
degrade somewhat after a point in certain settings.
